\documentclass[11pt,a4paper,sans]{moderncv}

% ----- MODERNCV SETUP -----
\moderncvstyle{classic}
\moderncvcolor{blue}

% ----- PACKAGES -----
\usepackage[utf8]{inputenc}
\usepackage[T1]{fontenc}
\usepackage[scale=0.84,top=1.2cm,bottom=1.2cm]{geometry}
\usepackage{lmodern}

% ----- PERSONAL DATA -----
\name{Luan Alfredo}{Lugarini}
\title{Computer Engineering Student}
\address{Curitiba}{Paraná, Brazil}
% \phone[mobile]{+55 (41) 9XXXX-XXXX}  % uncomment for the version sent by e-mail
\email{luanlugarini@gmail.com}
\social[github]{Luan1357}
\extrainfo{\href{http://lattes.cnpq.br/4564643932699945}{Lattes CV}}

\begin{document}

\makecvtitle

% ----- SUMMARY -----
\small{Fourth-semester Computer Engineering student at PUCPR seeking an engineering internship. Academic projects in systems programming in C (dynamic memory allocation, binary file parsing), Python and robotics (LEGO SPIKE and Arduino), including a robotic arm I assembled and programmed on my own. Advanced English (C1).}
\vspace{4pt}

% ----- EDUCATION -----
\section{Education}
\cventry{2025 -- 2029}{B.S. in Computer Engineering}{Pontifical Catholic University of Paraná (PUCPR)}{Curitiba, Brazil}{4th semester, expected graduation 2029}{}
\cventry{2021 -- 2023}{High School}{Colégio Imaculada Conceição}{}{}{}

% ----- PROJECTS -----
\section{Academic Projects}
\cventry{2025}{MyArm Robotic Arm: Tower of Hanoi}{Principles of Robotics, PUCPR}{C++ / Arduino}{}{
\begin{itemize}
    \item Assembled on my own a 4-degree-of-freedom robotic arm (base, two joints and gripper) driven by servo motors and an Arduino.
    \item Programmed the motion sequence that physically solves the Tower of Hanoi puzzle, with servo speed control for smooth movement.
    \item Code and video: \href{https://github.com/Luan1357/myarm-torre-de-hanoi}{github.com/Luan1357/myarm-torre-de-hanoi}
\end{itemize}}
\cventry{2025}{Autonomous Maze-Solving Robot}{Principles of Robotics, PUCPR}{Python / LEGO SPIKE}{}{
\begin{itemize}
    \item Robot that navigates a physical maze unassisted, using a distance sensor (walls), a color sensor (floor markers) and a gyroscope (90° and 180° turns with angle correction).
    \item Logic built as a state machine (outbound, return, finish) using asynchronous programming.
    \item \href{https://github.com/Luan1357/robo-labirinto-lego}{github.com/Luan1357/robo-labirinto-lego}
\end{itemize}}
\cventry{2026}{Conway's Game of Life}{Imperative Programming, PUCPR}{C}{}{
\begin{itemize}
    \item Cellular automaton simulator with a dynamically allocated grid (\texttt{malloc}/\texttt{calloc}), initial state loaded from file, and dead borders.
    \item Double buffering with pointer swapping between generations, avoiding data copies and incorrect in-place updates.
    \item \href{https://github.com/Luan1357/jogo-da-vida-conway}{github.com/Luan1357/jogo-da-vida-conway}
\end{itemize}}
\cventry{2026}{ppm2ascii: Image to ASCII Art Converter}{Imperative Programming, PUCPR}{C}{}{
\begin{itemize}
    \item Full parsing and validation of binary PPM (P6) images, supporting 8 and 16 bits per channel, with overflow checks before allocation.
    \item Per-pixel luminance calculation mapped onto a 57-character ramp.
    \item \href{https://github.com/Luan1357/ppm2ascii}{github.com/Luan1357/ppm2ascii}
\end{itemize}}
\cventry{2026}{15 Puzzle}{Imperative Programming, PUCPR}{Python}{}{
\begin{itemize}
    \item 4x4 sliding puzzle in the terminal, with a shuffle that always produces a solvable board and an auto-solver that replays the move history as a stack.
    \item \href{https://github.com/Luan1357/jogo-dos-15}{github.com/Luan1357/jogo-dos-15}
\end{itemize}}

% ----- SKILLS -----
\section{Technical Skills}
\cvitem{Languages}{C, Python, C++ (Arduino)}
\cvitem{Hardware}{Arduino, servo motors, LEGO SPIKE (distance, color and gyro sensors)}
\cvitem{Tools}{Git, GitHub, VS Code}

% ----- LANGUAGES -----
\section{Languages}
\cvitem{Portuguese}{Native}
\cvitem{English}{Advanced (C1), certified by Wise Up}

\end{document}
